\documentclass[11pt]{article}
\usepackage[margin=1in]{geometry}
\usepackage{amsmath,amssymb,amsthm}
\usepackage{booktabs,enumitem,microtype,xcolor}
\usepackage[hidelinks]{hyperref}
\usepackage{tcolorbox}
\newcommand{\figdir}{../figures}
% Shared colours and TikZ styles for the 3 x n Col papers.
\usepackage{tikz}
\usetikzlibrary{patterns,arrows.meta,positioning,decorations.pathreplacing,calc,fit,shapes.geometric}

\definecolor{colboth}{RGB}{255,255,255}
\definecolor{colblue}{RGB}{176,206,245}
\definecolor{colbluedark}{RGB}{28,86,176}
\definecolor{colwhite}{RGB}{252,214,160}
\definecolor{colwhitedark}{RGB}{214,128,20}
\definecolor{coldead}{RGB}{150,150,150}
\definecolor{case1}{RGB}{252,187,161}
\definecolor{case2}{RGB}{199,233,192}
\definecolor{case3}{RGB}{198,219,239}
\definecolor{case4}{RGB}{218,218,235}
\definecolor{case5}{RGB}{254,230,206}
\definecolor{case6}{RGB}{247,182,210}

\tikzset{
  gridline/.style={draw=black!30, line width=0.3pt},
  boardline/.style={draw=black!80, line width=0.7pt},
  blockline/.style={line width=1.1pt},
  blocklabel/.style={font=\scriptsize},
  boundlabel/.style={font=\tiny},
  seam/.style={draw=red!70!black, line width=1pt, densely dashed},
  bluestone/.style={circle, fill=colbluedark, draw=colbluedark!70!black, line width=0.4pt, inner sep=0pt},
  whitestone/.style={circle, fill=white, draw=black, line width=0.9pt, inner sep=0pt},
  basebox/.style={draw, rounded corners, fill=case2, font=\small, minimum width=1.5cm, minimum height=0.6cm},
  stepbox/.style={draw, rounded corners, fill=case3, font=\small, minimum width=1.5cm, minimum height=0.6cm},
  rootbox/.style={draw, rounded corners, fill=black!6, font=\small, inner sep=5pt},
  leafbox/.style={draw, rounded corners, fill=case3, font=\small, align=center, inner sep=4pt},
  smallbox/.style={draw, rounded corners, fill=case5, font=\scriptsize, align=center, inner sep=3pt},
}

\usepackage{adjustbox}
\newcommand{\colfig}[1]{\begin{adjustbox}{max width=\linewidth}\input{\figdir/gen/#1.tex}\end{adjustbox}}


\tcbuselibrary{skins}
\newtcolorbox{idea}[1][]{colback=colblue!25, colframe=colbluedark, fonttitle=\bfseries, title=#1}
\newtcolorbox{careful}[1][]{colback=red!5, colframe=red!60!black, fonttitle=\bfseries, title=#1}

\setlength{\parskip}{0.5em}
\setlength{\parindent}{0pt}

\title{\textbf{Who wins Col on a $3\times n$ board?}\\[4pt]
\large A friendly guide to a proof that the second player always wins}
\author{Gordon Chi}
\date{\today}

\begin{document}
\maketitle

\begin{idea}[The one-sentence answer]
On an empty board with 3 rows and any number of columns, the player who moves
\emph{second} can always win Col, no matter how cleverly the first player
plays.
\end{idea}

This guide explains why, step by step, with lots of pictures. You do not need
any background beyond knowing how to play a board game. A more formal version of the same
argument is written for mathematicians in the companion paper.

\tableofcontents

\section{Meet the game}

Col is a game for two players, whom we call \textbf{Blue} and \textbf{White}.
They take turns placing a stone of their own colour on an empty square of a
grid. There is exactly one rule:

\begin{idea}[The only rule]
You may \textbf{not} put a stone right next to (above, below, left or right
of) one of \textbf{your own} stones. Next to your opponent's stones is fine.
\end{idea}

If it is your turn and you have no legal square, you \textbf{lose}. There are
no draws, and there is no luck.

\begin{center}
\colfig{rules_forbidden}
\end{center}

As the game goes on, some squares become off-limits to one player, some to the
other, and some to both. We colour each empty square to show who may still
play there:

\begin{center}
\colfig{legend}
\end{center}

Here is a position on a board with 3 rows and 5 columns after four moves. The
numbers show the order the stones were played:

\begin{center}
\colfig{rules_legality}
\end{center}

Blue's stone \textbf{1} makes the squares around it pale orange: only White
may still play there. White's stone \textbf{2} makes its neighbours pale
blue. The grey squares are dead for both players. A player's legal squares
are the white ones plus those tinted in her own colour, so this picture shows
at a glance what each player can still do.

\section{What does ``the second player wins'' mean?}

Game theorists say a game has \textbf{value zero} if whoever moves second can
always win. It does not matter whether Blue or White starts; the one who waits
wins.

The smallest example is a strip of three squares. Here is every way the game
can go when Blue starts (the right end is a mirror image of the left, so we
skip it):

\begin{center}
\colfig{game_tree}
\end{center}

If Blue starts in the middle, White plays at an end and Blue is stuck, since
the other end is next to Blue's stone. If Blue starts at an end, White should
take the \emph{other} end. Then the middle is next to Blue's stone, and Blue is
stuck again. Either way, White wins. The same happens if White starts, with
the colours swapped. So this little strip has value zero.

\begin{idea}[Our goal]
Show that the empty board with 3 rows and $n$ columns (call it the
$3\times n$ board) has value zero, for \emph{every} $n$: 1, 2, 3, \dots,
1000, and forever.
\end{idea}

We can never try every $n$, and a single board can have more possible games
than there are atoms in the universe. We need ideas.

\section{Even widths: the copycat trick}

Suppose the board has an even number of columns, say 6. Imagine spinning the
board half a turn around its centre. Every square lands on a \emph{different}
square. The second player simply copies every move to the spun-around square:

\begin{center}
\colfig{mirror_even}
\end{center}

Why is the copy always legal? The board always looks the same after the spin,
with the colours swapped. If Blue's move was legal, then White's spun copy is
legal for exactly the same reason. So the copycat never runs out of moves, and
the other player gets stuck first.

\begin{careful}[Why this fails for odd widths]
With an odd number of columns, the centre square spins onto \emph{itself}. If
the first player takes the centre, there is nothing to copy onto.
\begin{center}\colfig{odd_center}\end{center}
The rest of this guide is about odd widths.
\end{careful}

\section{Two big ideas}

\subsection{Big idea 1: games can be added}

If a board splits into separate regions that do not touch, you can think of
it as several games played at once: on each turn you choose one region and
move there. Game theorists call this a \emph{sum} of games.

\begin{center}
\colfig{sum_intuition}
\end{center}

\begin{idea}[Adding second-player wins]
If White wins each region when Blue moves first there, White also wins the
whole sum. Her strategy: always answer in the same region Blue just played in.
\end{idea}

\subsection{Big idea 2: cheating in Blue's favour}

Our boards do not split into separate regions by themselves. So we pretend
they do. We cut the board into pieces and change the rules a little, always in
\textbf{Blue's} favour:
\begin{itemize}[nosep]
\item Blue may keep every square she really has, and we may give her extra
ones;
\item White may only use squares she really has, and we may take some away;
\item at each cut, White must not be allowed to play on both sides of the cut
in the same row.
\end{itemize}

\begin{center}
\colfig{comparison_schematic}
\end{center}

If White can still win this ``rigged'' game, she can certainly win the real
one. Her real moves copy her virtual moves, and Blue's real moves are also
available in the virtual game. This is the \textbf{comparison principle}, the
main tool of the whole proof.

\begin{careful}[Why the third rule matters]
In the real game, two squares side by side across a cut really are next to
each other. If White could play on both of them in the rigged game, she would
be allowed something the real rules forbid. So at every cut we check that
each row is White-legal on at most one side.
\end{careful}

\subsection{Pieces and plugs}

The pieces we cut the board into are called \textbf{gadgets}. Each gadget is a
small $3$-row board with some squares already restricted. The pattern of
White-legal squares on a gadget's left and right edges works like a
\emph{plug}. Two gadgets may sit side by side only if their plugs never both
have White in the same row:

\begin{center}
\colfig{seams}
\end{center}

Here is the whole toolbox. Under each gadget is a guarantee that a computer
has checked. ``$\le 0$'' means White wins that gadget when Blue moves first
there.

\begin{center}
\colfig{gadget_gallery}
\end{center}

\subsection{Scores that are not whole numbers}

Some gadgets come with a score like $-1$ or $\tfrac14$ instead of $0$. Think of
it as a count of spare moves: $-1$ means White is a whole spare move ahead, and
$+\frac14$ means Blue is a quarter of a move ahead. (Fractional moves really do
make sense in game theory!) When we add regions we add their scores. If the
total is below zero, White wins even if it is her turn to move.

\section{Widths $5, 9, 13, 17,\dots$ (one more than a multiple of 4)}

For these widths the pieces fit like a brick wall: repeat the four-column
gadget $E_4$ and finish with the one-column gadget $F_1$.

\begin{center}
\colfig{tiling_4k1}
\end{center}

Every plug fits, and every piece scores $\le0$. So White wins when Blue starts.
By the colour swap, Blue wins when White starts. So the board has value zero.

\begin{idea}[The colour-swap shortcut]
The empty board looks identical if you swap the names Blue and White. So if
the second player wins when Blue starts, the second player also wins when
White starts. We only ever need to check Blue starting.
\end{idea}

\section{Widths $3, 7, 11, 15,\dots$ (three more than a multiple of 4)}

This is the hardest family, and it needs a new idea: \textbf{induction}, the
domino effect. We show:
\begin{enumerate}[nosep]
\item the boards of width 3, 7 and 11 have value zero (checked by computer);
\item if every narrower board in the family has value zero, then so does the
next one.
\end{enumerate}
Then width 15 follows from 3, 7 and 11; width 19 from those and 15; and so on
forever, like falling dominoes.

\begin{center}
\colfig{recursion_chain}
\end{center}

\subsection{The starting dominoes}

For the $3\times 3$ board there are only four genuinely different first
moves for Blue. Here is White's winning answer to each:

\begin{center}
\colfig{base3}
\end{center}

Widths 7 and 11 are handled the same way, with 8 and 12 first moves, and a
computer checks every answer.

\subsection{Only a quarter of the board matters}

Flipping the board top-to-bottom or left-to-right does not change the game.
So we only need to answer Blue's first move in the shaded part:

\begin{center}
\colfig{normalization}
\end{center}

We sort those first moves into six cases, shown here for width 23:

\begin{center}
\colfig{case_map}
\end{center}

\subsection{The six cases, in pictures}

In each picture the top board is the real game after Blue's first move (and
White's answer, if there is one). The bottom board is the rigged version, cut
into gadgets. Circles show where the stones went.

\paragraph{Case 1: Blue plays on an ``odd'' square.}
Colour the board like a chessboard; ``odd'' squares are those where row plus
column is odd. Blue's move lands inside one of the $E_4$ bricks and damages it
so badly that it scores $-1$: White is a whole move ahead.

\begin{center}
\colfig{e4_parity}\qquad
\colfig{opened_e4}
\end{center}

The end-cap $C_3$ gives Blue back only a quarter:

\begin{center}
\colfig{number_line}
\end{center}

The total is $-\frac34$, below zero, so White wins even though it is her turn.
She does not need a specific answer at all!

\begin{center}
\colfig{example_case1}
\end{center}

\paragraph{Cases 2 to 5: special answers.}
When Blue plays an ``even'' square in the shaded area, White makes one specific
answer, and a special gadget around the two stones scores $\le 0$.

\begin{center}
\colfig{example_case2}\\[8pt]
\colfig{example_case3}\\[8pt]
\colfig{example_case4}\\[8pt]
\colfig{example_case5}
\end{center}

\paragraph{Case 6: the clever one.}
Blue plays in the middle row, in a column three more than a multiple of 4.
White answers diagonally up and to the right. Now look at Blue's column:

\begin{center}
\colfig{case6_zoom}
\end{center}

Blue can never play in that column again: one square has her stone, and the
other two are next to it. White can't use the top square, which is beside her
own stone, and she \emph{promises} never to use the bottom one. The column has
become a wall. To its left is an ordinary empty board, narrower than the one
we started with and in the same family, and we already know it has value zero
(the domino before). To its right sit certified gadgets.

\begin{center}
\colfig{example_case6}\\[8pt]
\colfig{example_case6_big}
\end{center}

\begin{idea}[Why this is the heart of the proof]
Case 6 is the only place the dominoes are used. One stone plus one answer
turns a big board into a smaller copy of the same problem. That is exactly
what induction needs.
\end{idea}

\section{The whole proof on one page}

\begin{center}
\colfig{roadmap}
\end{center}

\begin{itemize}
\item \textbf{Even width:} copy every move through the centre.
\item \textbf{Width $4k+1$:} a wall of $E_4$ bricks capped by $F_1$.
\item \textbf{Width $4k+3$:} three starting dominoes, then six cases for Blue's
first move. Case 1 wins on the score, Cases 2--5 use a special gadget, and
Case 6 builds a wall that leaves a smaller board.
\end{itemize}

\section{What did the computer do, and what did people do?}

Each gadget's guarantee concerns a small, finite board, so a computer can check
it completely. For each gadget the stored certificate lists, for every Blue
move, White's answer, and a short program replays all of these. Every
certificate is checked in a few seconds, without any searching.

The unlimited part, ``for every width'', cannot be checked by trying widths
one at a time. It rests on arguments: the copycat trick, the comparison
principle, the colour swap, the plug table, and the dominoes.

\begin{idea}[A computer checked the whole argument too]
All of those arguments have also been written in \textbf{Lean}, a
\emph{proof assistant}: a program that accepts a proof only if every single
step follows from the rules of logic. Lean checked the complete statement
``for every width $n$, the player who moves first on the empty $3\times n$
board loses'', stated with the real rules of Col, stones and all. It checked
the gadget certificates as well. Nothing is left to trust except Lean
itself.
\end{idea}

The Lean version uses an even simpler pattern for every case: bricks, then
one small \emph{window} around Blue's first move and White's answer, then
bricks again. The window is just the real board near the two stones, with
White giving up the few squares her plugs require.

\begin{center}
\colfig{lean_scheme}
\end{center}

Here are all nine windows it uses:

\begin{center}
\colfig{lean_windows}
\end{center}

With these windows the domino step already works at width 11, so the
starting dominoes are just the $3\times3$ and $3\times7$ boards.

\begin{careful}[What is still unknown]
Nobody yet knows whether the second player always wins on boards with 5 rows
and an odd number of columns, or on odd-by-odd boards in general. The wall
trick of Case 6 does not work there, because one stone does not block a whole
column of 5 squares.
\end{careful}

\section*{Glossary}
\begin{description}[nosep, font=\normalfont\bfseries]
\item[Value zero] The second player can always win.
\item[Gadget] A small certified piece of board with a score guarantee.
\item[Plug / port] The White-legal pattern on a gadget's left or right edge.
\item[Comparison principle] If White wins a version rigged in Blue's favour,
she wins the real game.
\item[Induction] Proving every case by proving the first ones and showing each
case follows from earlier ones.
\item[Certificate] A stored list of White's answers that a computer can replay
to confirm a gadget's guarantee.
\end{description}

\end{document}
